% Generated draft: requires mathematical and source review.
\documentclass[11pt]{article}
\usepackage[margin=25mm]{geometry}
\usepackage{fontspec}
\setmainfont{Latin Modern Roman}
\usepackage{amsmath,amssymb,mathtools,unicode-math}
\setmathfont{Latin Modern Math}
\usepackage{enumitem,needspace,xurl,hyperref}
\setlength{\parindent}{0pt}
\setlength{\parskip}{5pt}
\setlength{\emergencystretch}{3em}
\newcommand{\N}{\mathbb{N}}
\newcommand{\Z}{\mathbb{Z}}
\newcommand{\Q}{\mathbb{Q}}
\newcommand{\R}{\mathbb{R}}
\newcommand{\C}{\mathbb{C}}
\newcommand{\F}{\mathbb{F}}
\newcommand{\A}{\mathbb{A}}
\newcommand{\PP}{\mathbb{P}}
\newcommand{\E}{\mathbb{E}}
\newcommand{\eps}{\varepsilon}
\newcommand{\dd}{\,\mathrm{d}}
\begin{document}
\section*{UnsolvedMath ID: 20000531}
\textbf{Source problem code:} catalogue entry
\subsection{Definitions and mathematical statement}
Bennett's reconstructed AIM question concerns three families $\mathbb T_j$ of tubes in $\R^4$ with one side of length $1$ and three sides of length $\delta$, $0<\delta<1$. Directions within each family are $\delta$-separated, and uniform transversality means $\|e_1\wedge e_2\wedge e_3\|\geq c_0>0$ for every triple of family directions. Put $F_j=\sum_{T\in\mathbb T_j}\chi_T$ and $N_j=\#\mathbb T_j$. In the usual $1\leq p\leq\infty$ conjugate-exponent convention, $1/p+1/p'=1$, the question asks whether for every $\varepsilon>0$,
\[
 \|F_1F_2F_3\|_{L^{q/3}(\R^4)}\leq C_{\varepsilon,p,q,c_0}
 \prod_{j=1}^3\left[\delta^{4/q-3/p'-\varepsilon}N_j^{1/p}\right]
\]
whenever $q\geq4/3$ and $1/p+3/q\leq3$. The $L^{q/3}$ expression is a quasi-norm when $q<3$. The source does not fully state the $p$ range or quantitative transversality/multiplicity conventions; those matter for uniform constants.

For nonempty families relabeled $N_1\leq N_2\leq N_3$, the supplied report proves the unequal-cardinality estimate, for every loss parameter $\eta>0$,
\[
 \|F_1F_2F_3\|_{L^{13/27}}
 \leq C\delta^{99/13-\eta}N_1^{10/13}N_2N_3.
\]
At $(p,q)=(13/12,13/9)$ this is the conjectured cardinality product $(N_1N_2N_3)^{12/13}$ multiplied by
\[
 \left(\frac{N_2N_3}{N_1^2}\right)^{1/13}.
\]
The report obtains this by exact integer replication in the Guth--Zahl polynomial-Wolff estimate, with cancellation of the common replication scale against its $K^{1/9}$ dependence. The definition and verification of that replicated polynomial-Wolff constant are not supplied.

For $13/9\leq q\leq3/2$, let $1/p=3-3/q$ and $\theta=39(1/q-2/3)\in[0,1]$. Interpolation with the Bennett--Carbery--Tao epsilon endpoint is reported to give
\[
 \|F_1F_2F_3\|_{L^{q/3}}\leq C\delta^{3(4/q-3/p')-\eta}
 \prod_jN_j^{1/p}\left(\frac{N_2N_3}{N_1^2}\right)^{\theta/13}.
\]
Losses $\eta$ and the source's three factorwise epsilon losses are equivalent by renaming an arbitrary positive parameter. For comparable family sizes, with comparison constants independent of $\delta$, the imbalance is bounded; the report consequently obtains the requested admissible range $q\geq13/9$ for such families. It does not reach the whole conjectured $4/3\leq q<13/9$ range or remove the unequal-size loss.

Finally the report claims two affine-foliation constructions rigorously recover both necessary conditions $q\geq4/3$ and $1/p+3/q\leq3$. Their details and full theorem references are absent; necessity is not a sufficiency proof.
\subsection{Short English statement}
Seek the trilinear direction-separated Kakeya bound in four dimensions throughout its stated exponent range. The report gives a specific endpoint with an explicit family-size imbalance loss and interpolates it. Comparable sizes recover a smaller admissible range; general unequal sizes and the remaining conjectured range are not settled.
\subsection{Sources}
\begin{itemize}
\item Local registry: \url{data/unsolvedmath/problems.json}, permanent ID $20000531$, source code \texttt{catalogue entry}; damaged original equation (0.3), supplied reconstruction and research summary. Statement status reconstructed and unverified.
\item AIM workshop ``Carleson theorems and multilinear operators'', Question 4, attributed to Bennett: \url{https://aimath.org/pastworkshops/multilinopsproblems.pdf}. Supplied verification identifies page 2; canonical location \texttt{aim-analysis-notes.json}, \texttt{notes[62]}. No PDF was opened here.
\item The summary names Guth--Zahl's polynomial-Wolff estimate and Bennett--Carbery--Tao's epsilon endpoint without full citations, precise hypotheses or theorem numbers. Original literature field empty.
\item Supplied machine-generated research attempt 1: partial result; problem partially solved at the run and in the registry; rigor metadata ``proved'', novelty confidence low. The non-peer-reviewed replication/interpolation and necessary-condition claims require independent verification before citation.
\end{itemize}
\end{document}
